\originalsection{作业3}\label{sec:ps03}
原件3页, 数学题面2页.
\begin{exercise}\label{ps03:1}
原题1. 令 $M=\{a\in\ell^\infty:\lim a_k\text{存在}\}$. (a) 证明 $u(a)=\lim a_k$ 是 $M$ 上的有界线性泛函. (b) Hahn--Banach 给出延拓 $v\in(\ell^\infty)^*$, $v|_M=u$; 证明不存在 $b\in\ell^1$ 使 $v(a)=\sum a_kb_k$ 对所有 $a\in\ell^\infty$ 成立. 原提示: 测试 $e_n$.
\end{exercise}
\begin{proof}[AI 独立解答]
(a) 极限的线性性给出线性, $|u(a)|\le\norm a_\infty$ 给出有界. 常数数列1说明 $\norm u=1$. (b) 因 $e_n\in M$ 且极限为零, 若存在此 $b$, 则 $b_n=v(e_n)=0$ 对所有 $n$ 成立. 但 $v(1,1,\ldots)=1$, 与零配对矛盾. 这复用第\ref{ch:dual}章关于 $\ell^\infty$ 对偶的例题解答, 未对 $p=\infty$ 套用题\ref{ps01:5}的结论.
\end{proof}
\begin{exercise}\label{ps03:2}
原题2. 设 $V,W$ 为赋范空间, $T\in\B(V,W)$. (a) 定义 $T^\dagger f=f\circ T$ 于 $f\in W^*$; 证明 $T^\dagger\in\B(W^*,V^*)$ 且 $\norm{T^\dagger}=\norm T$. (b) 对 $1\le p\le\infty$, 证明右移 $Ra=(0,a_1,a_2,\ldots)$ 在 $\ell^p$ 上有界, 并求范数. (c) 对 $1\le p<\infty$ 及共轭指数 $q$, 用题\ref{ps01:5}的配对识别对偶为 $\ell^q$, 求 $R^\dagger b$. 原提示(a): Hahn--Banach 给出达到 $\norm{Tx}$ 的范数1泛函; (c)以 $R^\dagger e_1=0$ 为例.
\end{exercise}
\begin{proof}[AI 独立解答]
(a) 对每个 $f$, $f\circ T$ 线性且 $\norm{f\circ T}\le\norm f\norm T$; 关于 $f$ 也线性, 故得上界. 若 $Tx\ne0$, 推论\ref{cor:norming}给出 $\norm f=1$, $f(Tx)=\norm{Tx}$, 所以 $\norm{Tx}\le\norm{T^\dagger}\norm x$. $Tx=0$ 时同式直接成立. 对单位向量取上确界得反向界, 零空间亦满足等式.

(b) 右移不改变有限 $p$ 的幂和或上确界, 所以 $\norm{Ra}_p=\norm a_p$, $\norm R=1$.

(c) 绝对收敛允许重编号: $(R^\dagger b)(a)=\sum_{k\ge2}a_{k-1}b_k=\sum_{j\ge1}a_jb_{j+1}$. 因而 $R^\dagger b=(b_2,b_3,\ldots)$, 是 $\ell^q$ 上的左移; 此也包括 $p=1,q=\infty$.
\end{proof}
\begin{exercise}\label{ps03:3}
原题3. 对 $E\subset\R$, $x\in\R$, 令 $E+x=\{y+x:y\in E\}$, 证明 $m^*(E+x)=m^*(E)$.
\end{exercise}
\begin{proof}[AI 独立解答]
任意覆盖 $E$ 的开区间列 $(I_n)$ 平移为覆盖 $E+x$ 的开区间列 $(I_n+x)$, 长度总和不变. 对覆盖取下确界得 $m^*(E+x)\le m^*(E)$. 再以 $-x$ 平移 $E+x$ 得反向界. 两边为无穷时论证仍成立, 不涉及无穷相减.
\end{proof}
\begin{exercise}\label{ps03:4}
原题4. 设 $U\subset\R$ 为非空开集, 对 $x\in U$ 令 $a_x=\inf\{a\in\R:(a,x]\subset U\}$, $b_x=\sup\{b\in\R:[x,b)\subset U\}$. (a) 证明 $(a_x,b_x)\subset U$. (b) 若 $y\in(a_x,b_x)$, 证明 $(a_y,b_y)=(a_x,b_x)$. (c) 证明 $U=\bigcup_{q\in U\cap\Q}(a_q,b_q)$. 原提示: $\Q$ 在 $\R$ 中稠密. 端点允许取 $\pm\infty$; 这是原定义在无界开集上的必要解释.
\end{exercise}
\begin{proof}[AI 独立解答]
(a) 开性给出 $a_x<x<b_x$. 若 $a_x<t<x$, 下确界定义给出某个 $a<t$ 使 $(a,x]\subset U$, 所以 $t\in U$; 右侧同理, $x$ 自身也在 $U$.

(b) 区间 $I_x=(a_x,b_x)$ 是包含 $x$ 且包含于 $U$ 的最大开区间: 若另一个这样的区间左端更小, 其与 $I_x$ 的并仍为 $U$ 内包含 $x$ 的开区间, 与 $a_x$ 的定义矛盾; 右端同理. 对 $y\in I_x$, $I_y$ 与 $I_x$ 相交, 二者的并为 $U$ 内开区间, 含 $x$ 和 $y$. 分别用两个最大性得 $I_x\cup I_y=I_x=I_y$.

(c) 每个 $I_q\subset U$. 反之, 对 $x\in U$ 在 $I_x$ 内选有理数 $q$, 由(b)得 $I_q=I_x$, 所以 $x$ 属于右端. 指标集可数, 给出可数个开区间的并; 重复区间可去重, 不同区间由(b)必不交.
\end{proof}
